\chapter{Relative Krajewski star: вывод support module и двухтекстурный no-go}

Предыдущий гейт условно потребовал, чтобы каждый subleading connector касался
выбранного state \(\rho_\star\). Теперь проверим, является ли это отдельным
postulate или следует из минимального finite graph после выбора state.

\section{Классификация corner bimodules}

Для \(Q_\star=I-\rho_\star\) полная family matrix algebra раскладывается в
четыре \(\mathcal A_\rho\)-bimodule corners:
\begin{equation}
 M_3=
 \rho M_3\rho
 \oplus\rho M_3Q
 \oplus QM_3\rho
 \oplus QM_3Q,
\end{equation}
с комплексными размерностями
\begin{equation}
 1,\qquad2,\qquad2,\qquad4.
\end{equation}
Полный перебор всех шестнадцати сумм corners с условием *-замкнутости оставляет
восемь candidates. Наложим три graph requirements:
\begin{enumerate}
\item anchor self-edge \(\rho M_3\rho\) сохраняет leading heavy direction;
\item edges \(\rho M_3Q\) и \(QM_3\rho\) связывают anchor с двумя light
directions;
\item внутренний complement edge \(QM_3Q\) отсутствует как не связанный с
выбранным boundary anchor.
\end{enumerate}
Этим требованиям удовлетворяет ровно один *-замкнутый bimodule:
\begin{equation}
 \rho M_3\rho\oplus\rho M_3Q\oplus QM_3\rho
 =\mathcal M_\rho,
 \qquad \dim\mathcal M_\rho=5.
\end{equation}

\begin{theorem}[Unique relative-star module]
Среди *-замкнутых \(\mathcal A_\rho\)-bimodules единственный минимальный graph,
который содержит anchor loop, соединяет anchor с complement и не вводит
самостоятельную light--light dynamics, есть пятимерный relative Krajewski
star \(\mathcal M_\rho\).
\end{theorem}

В basis, где \(\rho=\operatorname{diag}(1,0,0)\), его пять matrix units ---
ровно entries первой строки или первого столбца. Возвращение к произвольному
basis даёт прежнее covariant условие \(QXQ=0\).

\section{Projection как endpoint inclusion--exclusion}

Каждый исходный incidence operator \(H\) имеет два anchored restrictions:
\(\rho H\) и \(H\rho\). Их пересечение \(\rho H\rho\) посчитано дважды.
Обычная inclusion--exclusion поэтому даёт
\begin{equation}
 \rho H+H\rho-\rho H\rho
 =H-QHQ
 =\Pi_\rho(H).
\end{equation}
Следовательно, коэффициенты \((1,1,-1)\) projection не являются flavour
fit: это целочисленная endpoint incidence формула relative star.

\begin{proposition}[Conditional derivation of the support axiom]
Если subleading family graph определяется как минимальный relative graph к
вариационно выбранному anchor, то support axiom и projection \(\Pi_\rho\)
следуют из incidence и *-замкнутости, а не вводятся независимо.
\end{proposition}

\section{Почему четырёх Yukawa-текстур ещё нет}

Observed representation различает четыре edge types, но family graph пока
факторизован от gauge labels. Поэтому
\begin{equation}
 u,\nu\longmapsto H_{\rm up},
 \qquad
 d,e\longmapsto H_{\rm down}.
\end{equation}
Получаются только две normalized family textures. Цветовая multiplicity
умножает quark functional на \(3\), но положительное scalar умножение не
меняет eigenvectors и variational ground state. Exact test даёт ground-state
residual меньше \(2\times10^{-16}\).

\begin{theorem}[Factorized four-sector no-go]
При gauge--family factorization центральные quark/lepton projectors и их
multiplicities не расщепляют family state. Один relative star даёт максимум
две textures, связанные с \(H\) и \(\widetilde H\), но не четыре независимых
operators \(Y_u,Y_d,Y_e,Y_\nu\).
\end{theorem}

\begin{nogocriterion}[Запрет четырёх ручных anchors]
Нельзя независимо выбрать \(\rho_u,\rho_d,\rho_e,\rho_\nu\). Следующий pass
допустим только если quark--lepton splitting возникает из одного общего
representation-sensitive functional и сохраняет anomaly/order-one ledger.
\end{nogocriterion}

Итак, support module выведен условно из минимального relative graph. Новый
узкий gap --- получить неfactorizable quark/lepton dependence, не возвращая
четыре свободные Yukawa matrices. Машинный classification ledger сохранён в
\path{s2t_v4_relative_krajewski_star_gate_results.json}.